\section{Conclusion and Future Work}
\label{conclusion}
This study presented a low-cost, machine vision-based virtual encoder for 3D printer kinematic calibration, integrating a planar chessboard marker, an EDLines extraction pipeline, and an iterative SolvePnP solver. Characterization testing established a mean reprojection error of $0.1192 \pm 0.0021\text{ px}$, an internal static noise floor of $2.58~\mu\text{m}$ ($\sigma = 2.5~\mu\text{m}$) on the X-axis and $4.9~\mu\text{m}$ ($\sigma = 3.3~\mu\text{m}$) on the Y-axis, and absolute measurement errors of $3.55~\mu\text{m}$ and $6.25~\mu\text{m}$, respectively. Following fine-tuning, the calibration framework systematically mitigated all linear error components: X-axis positioning, skew, and backlash errors were reduced by 84.6\% (to $7.1 \pm 5.9~\mu\text{m}$), 93.3\% (to $6.3 \pm 3.7~\mu\text{m}$), and 45.6\% (to $6.8 \pm 6.8~\mu\text{m}$), while Y-axis metrics decreased by 90.2\% (to $5.4 \pm 5.4~\mu\text{m}$), 94.5\% (to $5.0 \pm 4.3~\mu\text{m}$), and 72.0\% (to $6.8 \pm 5.2~\mu\text{m}$). Furthermore, dynamic trajectory tracking error was suppressed by 54.24\%, yielding a total path RMSE of $46.17 \pm 1.77~\mu\text{m}$. These empirical findings confirm that software-driven error compensation effectively resolves multi-axis mechanical imperfections under a hardware budget below 30 USD. 

Future work will focus on integrating telecentric optics to minimize perspective distortion and developing a rigid, CNC-machined enclosure mounted on a modular toolhead to reduce structural flexing during dynamic calibration.